\section*{\centering \sffamily Prólogo}

Las VIII Jornadas de Enseñanza de la Matemática (JEM), realizadas en 2024, constituyeron un espacio de encuentro, reflexión e intercambio entre docentes, investigadores, estudiantes y profesionales interesados en la enseñanza y el aprendizaje de la matemática en todos los niveles educativos.

En esta edición, marcada por el compromiso de fortalecer la calidad educativa, se compartieron experiencias, investigaciones y propuestas innovadoras que enriquecieron la mirada sobre la didáctica de la matemática, el uso de tecnologías y la diversidad de enfoques para potenciar el aprendizaje.

Las actividades incluyeron conferencias magistrales, mesas redondas, comunicaciones breves y talleres, en los que se abordaron temáticas actuales y relevantes para la comunidad educativa. Este variado programa permitió fortalecer los lazos profesionales y académicos, a la vez que abrió nuevas oportunidades de colaboración e investigación.

A través de estas Memorias, ponemos a disposición de la comunidad los trabajos presentados durante las jornadas, con el propósito de difundir el conocimiento generado y favorecer su aplicación en diversos contextos.

El Comité Organizador agradece profundamente a las y los participantes, disertantes, moderadores, instituciones y colaboradores que hicieron posible la realización de estas jornadas. Su compromiso y dedicación son la base para seguir construyendo un espacio abierto, inclusivo y de permanente crecimiento para la enseñanza de la matemática.

Cerramos esta edición con la convicción de que el intercambio y la cooperación seguirán siendo motores fundamentales para mejorar nuestras prácticas y, con ello, la calidad educativa.

\vspace{5cm}

\begin{center}
	{\Large Comité Organizador}
	
	{\normalfont VIII Jornadas de Enseñanza de la Matemática – 2024}
\end{center}